% © 2026 Ing. David Jaroš — CC BY-NC-ND 4.0
% UBT-AI-PROVENANCE-BEGIN
% schema: ubt-ai-provenance/v1
% tier: B_machine_verified
% ai_assistance: disclosed
% human_review: machine-verification
% editorial_responsibility: Ing. David Jaroš
% policy: ../../AI_PROVENANCE.md
% notice: Elementary convention identities are exact; standard theta transformation/product theorems are cited as external mathematics.
% UBT-AI-PROVENANCE-END

\documentclass[11pt,a4paper]{article}
\usepackage[margin=2.5cm]{geometry}
\usepackage{amsmath,amssymb,amsthm}
\usepackage[most]{tcolorbox}
\newtheorem{theorem}{Theorem}
\newtheorem{corollary}[theorem]{Corollary}
\newtheorem{remark}[theorem]{Remark}
\title{Theta and Modular Variables in UBT:\\
Canonical Separation, Congruence-Subgroup Structure, and No-Go Results}
\author{Ing.\ David Jaro\v{s}}
\date{9 October 2026}
\begin{document}
\maketitle

\section{Three variables that must not be conflated}

\paragraph{Canonical UBT complex time.}
\[
 \tau_{\rm UBT}=t+i\psi,
 \qquad \psi\sim\psi+2\pi R_\psi .
\]
With noncompact physical real time, the coordinate domain is locally
\[
 \mathbb R_t\times S^1_\psi,
\]
a cylinder.  A single compact circle does not by itself define a two-torus
modulus.

\paragraph{Heat-kernel Jacobi modulus.}
The exact compact-fibre heat trace is
\[
 K_\psi(s)=\sum_{n\in\mathbb Z}e^{-s n^2/R_\psi^2}
 =\vartheta_3(0|\tau_H),
 \qquad
 \boxed{\tau_H=\frac{i s}{\pi R_\psi^2}}.
\]
Here $s>0$ is Euclidean proper time.  This is an exact spectral parameter,
not physical complex time.

\paragraph{Conditional geometric/thermal torus modulus.}
If an additional Euclidean real-time periodicity
$t_E\sim t_E+\beta$ is imposed, the two periods
$\beta$ and $2\pi R_\psi$ define a rectangular torus with a dimensionless
shape modulus proportional to
\[
 \tau_{\rm th}=i\frac{2\pi R_\psi}{\beta}
\]
up to the chosen cycle normalization.  This is a conditional thermal/geometric
construction.  It is not the coordinate $t+i\psi$ itself.

\section{Theta convention audit}

Use the standard Jacobi convention
\[
 \vartheta_3(0|\tau)
 =\sum_{n\in\mathbb Z}e^{\pi i\tau n^2},
 \qquad \operatorname{Im}\tau>0.
\]
Then the standard transformations include
\[
 \boxed{\vartheta_3(0|\tau+1)=\vartheta_4(0|\tau),}
\]
\[
 \boxed{\vartheta_3(0|\tau+2)=\vartheta_3(0|\tau),}
\]
and
\[
 \boxed{
 \vartheta_3(0|-1/\tau)
 =(-i\tau)^{1/2}\vartheta_3(0|\tau)
 }.
\]

Therefore a statement that the standard $\vartheta_3$ constant is invariant
under the generator $T:\tau\mapsto\tau+1$ is false.

A second convention frequently used in older UBT notes is
\[
 q=e^{2\pi i\tau},
 \qquad
 \Theta_{\mathbb Z}(\tau)=\sum_n q^{n^2}.
\]
But then
\[
 \boxed{
 \Theta_{\mathbb Z}(\tau)=\vartheta_3(0|2\tau).
 }
\]
Its $T$-invariance follows because the Jacobi argument changes by two.
One may not combine that $T$ law with the full-$S$ formula for
$\vartheta_3(0|\tau)$ without also transforming the factor of two.

The lattice theta series
\[
 \Theta_{\mathbb Z}(\tau)=\sum_n e^{2\pi i\tau n^2}
\]
is naturally a weight-$1/2$ theta modular form for a congruence subgroup
(with the usual half-integral-weight multiplier structure), not an ordinary
scalar modular form for the full $SL(2,\mathbb Z)$ with the naive pair of
$S,T$ laws.  Its $d$-fold product has weight $d/2$ only in the corresponding
theta-series modular setting.

For the current minimal $S^1_\psi$ spectrum this gives $d=1$ and weight
$1/2$ at the arithmetic theta-series level.  The historically used
$\vartheta_3^3$ and weight $3/2$ require a three-dimensional compact lattice,
which is not part of the minimal canonical compact fibre.

\section{No theta-3 zero in the upper half-plane}

Jacobi's product formula gives, for $|q|<1$,
\[
 \vartheta_3(0,q)
 =\prod_{n=1}^{\infty}
 (1-q^{2n})(1+q^{2n-1})^2 .
\]
Every factor is nonzero for $|q|<1$, and the convergent product is nonzero.
Hence
\[
 \boxed{
 \vartheta_3(0|\tau)\ne0
 \quad\text{for all }\tau\in\mathbb H.
 }
\]

Therefore any UBT mechanism that assumes a zero of
$\vartheta_3(0|\tau)$ at $(1+i)/2$ or elsewhere in the upper half-plane is
closed as a mathematical no-go.

\section{What remains genuinely modular in UBT}

The strongest exact modular statement in the current minimal theory is the
Poisson/Jacobi duality of the compact-$\psi$ heat trace:
\[
 \sum_n e^{-s n^2/R_\psi^2}
 =
 \sqrt{\frac{\pi R_\psi^2}{s}}
 \sum_k e^{-\pi^2R_\psi^2k^2/s}.
\]
This is exactly the Jacobi $S$ transformation evaluated on the spectral
modulus $\tau_H=i s/(\pi R_\psi^2)$.

It proves a dual representation of the same KK spectral trace.  It does
\emph{not} by itself prove:
\begin{itemize}
 \item modular invariance of the full UBT action;
 \item a modular action on physical complex time $\tau_{\rm UBT}$;
 \item a self-dual dynamical vacuum;
 \item a Hecke action on physical particle states;
 \item a Kac--Moody level or the exponent in an effective coupling.
\end{itemize}

\section{P1 classification status}

\begin{tcolorbox}[colback=green!6!white,colframe=green!55!black,title=Closed]
\begin{itemize}
 \item GAP-THETA-TAU-NAIVE: closed as a no-go.
 \item GAP-THETA-KK-MODULUS: closed in the conditional spectral branch;
       $\tau_H=i s/(\pi R_\psi^2)$.
 \item GAP-THETA-T-LAW: closed; convention mismatch explains the historical
       contradictory $T$ transformations.
 \item GAP-THETA3-ZERO: closed as a no-go; the theta constant has no zeros in
       the upper half-plane.
 \item GAP-COMPACT-MINIMAL: the minimal compact internal fibre is one circle,
       so the canonical heat trace is one-dimensional.
\end{itemize}
\end{tcolorbox}

\begin{tcolorbox}[colback=yellow!8!white,colframe=orange!70!black,title=Open]
A physical modular symmetry of the full UBT action remains open.  It requires
an action-derived dimensionless modulus and a transformation law for the full
field, measure, interactions and constrained fluctuation space.  Mock-modular
or resurgent extensions should be investigated only for an actual UBT kernel
that fails ordinary modularity in the corresponding way.
\end{tcolorbox}

\section{Standard references}
The transformation and product identities used here are standard Jacobi theta
identities; see NIST DLMF, Chapter 20 (Theta Functions), especially the
definitions, transformation identities and product formula for
$\vartheta_3(0,q)$.

\end{document}
